% !TEX root = approx8.tex

\section{Triangulated persistence refinements of Fukaya
  categories} \label{sec:alg1}
The aim of this section is to fix the basic techniques and notation to
deal with triangulated persistence categories and their applications
to Fukaya categories. Persistence categories and some of their
properties are discussed in \S\ref{subsec:PC}. One important feature
mentioned here is that the set of objects of such a category carries a
class of natural pseudo-distances called {\em interleaving}, see
\eqref{eq:d-int}. Triangulated persistence categories are briefly
recalled in \S\ref{subsec:TPC_approx} where is also introduced TPC
approximability in Definition \ref{def:TPC-approx}. This is the
version of approximability appearing in our main theorem, Theorem
\ref{thmmain1}. In \S\ref{subsec:A_{infty}} we start to discuss
filtered $A_{\infty}$ categories. The key section in this respect is
\S\ref{sb:pdc} where are discussed derived categories associated to
filtered $A_{\infty}$ categories. Section \ref{s:sys-tpc} makes
explicit the formal properties of a system of $A_{\infty}$ categories
(and the associated derived categories) that one obtains by making
choices of perturbations that become smaller and smaller. In
\S\ref{s:approximability} a new, more precise, version of
approximability is defined taking into account the system of
categories with increasing precision discussed before.  This is a
stronger, but more technical, version of Definition
\ref{def:TPC-approx} which is shown later in the paper to be satisfied
in the cases of geometric interest. Finally, \S\ref{s:pdfuk} contains
the main steps in the construction of the filtered Fukaya categories
to which can be applied the algebraic tools described earlier in the
section. We refer to \cite{BCZ:tpc} for further details on the
filtered algebra background material in this section.


\subsection{Persistence categories}\label{subsec:PC}
In brief, a persistence category $\mathscr{C}$ is a category enriched
over the category of persistence modules.  We fix notations and give the main
definitions below.
\subsubsection{Persistence modules}
A persistence module
$$M = \Bigl( \{M^{\alpha}\}_{\alpha \in \mathbb{R}}, \; i_{\alpha,
  \beta}: M^{\alpha} \longrightarrow M^{\beta}, \; \forall \ \alpha
\leq \beta \Bigr)$$ over a ring $R$ (for us $R$ will be often
  either a field or the positive Novikov ring) is a collection of
$R$-modules $M^{\alpha}$ for each $\alpha\in\R$ and
morphisms $i_{\alpha,\beta}$ such that
$i_{\beta,\gamma}\circ i_{\alpha,\beta}=i_{\alpha,\gamma}$ whenever
$\alpha\leq\beta\leq\gamma$ and $i_{\alpha,\alpha}=id$ (sometimes
additional constraints are imposed on these modules, as needed).  Let $u, v \in M^{\alpha}$ and $r \geq 0$. We write
  $u \underset{r}{=}v$ to indicate that $i_{\alpha,\alpha+r}(u) = i_{\alpha,\alpha+r}(v)$.
%\item For $\alpha \leq \beta$ we write
%  $m_{\beta}: M^{\alpha} \longrightarrow M^{\beta}$ for the map
%  $i_{\alpha, \beta}$.
%\item For $u \in M^{\alpha}$, $v \in M^{\beta}$. We write
%  $u \underset{\text{per}}{=} v$ to indicate that
%  $m_{\gamma}(u) = m_{\gamma}(v)$, where
%  $\gamma := \max \{\alpha, \beta \}$.
% For $u \in M^{\alpha}$ we denote 
%  $\ell(u) := \alpha$ and call $\ell(u)$ the persistence level of $u$.
%
We  denote by
$$M_{\text{tot}} := \coprod_{\alpha \in \mathbb{R}} M^{\alpha}$$ the underlying set
of $M$ %Note that the persistence level $\ell$ can be viewed as a function $\ell: M_{\text{tot}} \longrightarrow \mathbb{R}$.  
and we define $$M^{\infty} := \colim_{\alpha \to \infty} M^{\alpha},$$
where the colimit (or direct limit) is taken with respect to the
structural maps $i_{\alpha, \beta}$, $\alpha \leq \beta$, and call it
the $\infty$-limit of $M$.

\begin{rem}
  Let $C$ be a filtered chain complex, with an increasing filtration
  $C^{\alpha} \subset C^{\beta} \subset C$,
  $\forall \; \alpha \leq \beta$. We will sometimes view $C$ as a
  persistence module with
  $i_{\alpha, \beta}: C^{\alpha} \longrightarrow C^{\beta}$ being the
  inclusion maps. Let $u \in C^{\alpha}$ and consider the same
  element $u$ but now viewed in $C^{\beta}$ for some $\beta > \alpha$.
  For many practical purposes these two elements are viewed as one
  element that belongs to $C$. However, from the persistence module
  viewpoint $u \in C^{\alpha}$ and its image
  $i_{\alpha, \beta}(u) \in C^{\beta}$ should be viewed as distinct
  elements. \end{rem}


\subsubsection{Persistence categories} \label{sb:pc} (PC's in
short) are (small) categories in which the morphisms between objects form
persistence modules and the composition of morphisms is compatible
with the persistence structures. The general theory of persistence
categories is described in detail in~\cite[Section~2]{BCZ:tpc} and
here we will only recall several key concepts.

Unless otherwise stated, from now on we will always assume all the PC
categories to be endowed with a shift functor (or, more precisely, a
system of shift functors) $\Sigma = \{\Sigma^r\}_{r \in \mathbb{R}}$
(see~\cite[Section~2.2.3]{BCZ:tpc} for the definition), and all the PC
functors to commute with $\Sigma^r$ for every $r \in \mathbb{R}$. Part
of the shift functor structure are the natural transformations
  $\eta_{r,s}: \Sigma^r \longrightarrow \Sigma^s$,
  $s, r \in \mathbb{R}$.  The special case $s=0$ will be especially
  important and we denote it by
  $\eta_r := \eta_{r,0}: \Sigma^r \longrightarrow \id$.
Moreover, all persistence functors
  $\msf: \msc \longrightarrow \msd$ between PC's will be implicitly
  assumed to be compatible with the shift functors of these categories
  in the following sense: $\msf \circ \Sigma^r = \Sigma^r \circ \msf$
  for all $r \in \mathbb{R}$, and
  $(\eta_{r,s})_{\msf A} = \msf ((\eta_{r,s})_A)$ for all
  $A \in \Ob(\msc)$ and $r, s \in \mathbb{R}$.

Let $\mathscr{C}$ be a PC category. Denote by $\mathscr{C}^0$ the
$0$-level category associated with $\mathscr{C}$ and by
$\mathscr{C}^{\infty}$ the limit (or $\infty$-level) category of
$\mathscr{C}$. Isomorphisms in $\msc^0$ will be called 
$0$-isomorphisms and the property of being $0$-isomorphic will be 
sometimes denoted by $\cong_0$. Similarly, the existence of an isomorphism 
in $\msc^{\infty}$ will be denoted by $\cong_{\infty}$.

Most of the time we will assume our PC's $\mathscr{C}$ to be graded in
the sense that the persistence modules $\hom_{\mathscr{C}}(A,B)$ are
$\mathbb{Z}$-graded. We will mostly use cohomological grading
conventions and denote by $\hom_{\mathscr{C}}(A,B)^k$,
$k \in \mathbb{Z}$ the degree-$k$ component of
$\hom_{\mathscr{C}}(A,B)$. For $\alpha \in \mathbb{R}$ we denote by
$\hom^{\alpha}_{\mathscr{C}}(A,B)^k$ the $\alpha$-persistence level of
$\hom_{\mathscr{C}}(A,B)^k$. We assume composition of morphisms
to be degree-preserving and that identity morphisms are in degree $0$.

Our PC's $\mathscr{C}$ will be often endowed with a {\em translation
  functor} $T: \mathscr{C} \longrightarrow \mathscr{C}$. This is a PC
isomorphism functor with the property that for all
$A, B \in \Ob(\mathscr{C})$, $k \in \mathbb{Z}$, we have PC-natural
isomorphisms
$$\hom_{\mathscr{C}}(TA,B)^k \cong \hom_{\mathscr{C}}(A,B)^{k-1},
\quad \hom_{\mathscr{C}}(A,TB)^k \cong
\hom_{\mathscr{C}}(A,B)^{k+1}.$$ As mentioned earlier, since $T$ is a
PC functor, we will implicitly assume that $T$ commutes with the shift
functors, i.e.~$T \circ \Sigma^{r} = \Sigma^r \circ T$,
$\forall \, r \in \mathbb{R}$.




\subsubsection{Acyclic objects} \label{sbsb:acyclic} Let $\mathscr{C}$
be a PC, $A \in \Ob(\mathscr{C})$ and $r \geq 0$. We say that $A$ is
$r$-acyclic if the morphism $\eta_r^A: \Sigma^r A \longrightarrow A$
equals $0$, or, equivalently, if the identity $id_{A}$ satifies
$i_{0,r}(id_{A})=0$.  Here,
$\eta^A_{r} \in \hom_{\mathscr{C}^0}(\Sigma^{r}A, A)$ is given by
$\eta^{A}_{r}=\eta_{r}(A)$. An object $A$ is called acyclic if it is
$r$-acyclic for some $r \geq 0$. Equivalently, $A$ is acyclic if and
only if $A$ is isomorphic to $0$ in the limit category
$\mathscr{C}^{\infty}$.  We denote by $A\mathscr{C}$ the full
subcategory of $\mathscr{C}$ whose objects are the acyclic ones.


\subsubsection{Pseudo metrics associated with
  PC's} \label{sbsb:pc-metrics}

Let $\mathscr{C}$ be a PC with a shift functor $\Sigma$. The
collection of objects $\Ob(\mathscr{C})$ carries a natural pseudo-metric
$\dint$, called the interleaving distance, which is defined as
follows.  For $X, Y \in \Ob(\mathscr{C})$ define:
\begin{equation} \label{eq:d-int}
  \begin{aligned}
    \dint(X, Y) = \inf \Bigl\{r \geq 0 & \bigm| \exists \, \varphi:
    \Sigma^rX \longrightarrow Y, \, \exists \, \psi: \Sigma^rY
    \longrightarrow X \\
    & \;\;\; \text{such that} \; \psi \circ \Sigma^r \varphi =
    \eta_{2r}^X, \; \varphi \circ \Sigma^r \psi = \eta_{2r}^Y \Bigr\}.
  \end{aligned}
\end{equation}
Here, $\eta^X_{2r} \in \hom_{\mathscr{C}^0}(\Sigma^{2r}X, X)$ are the
standard maps associated with the persistence structure of
$\mathscr{C}$ and its shift functor and similarly for $\eta^Y_{2r}$.

Note that $\dint$ is in general not a genuine metric but only a
pseudo-metric and moreover it may have infinite values. Indeed, if $X$
and $Y$ are isomorphic in the subcategory $\mathscr{C}^0$, then
$\dint(X, Y) = 0$. Similarly, if $X$ and $Y$ are not isomorphic in
$\mathscr{C}^{\infty}$ then $\dint(X,Y) = \infty$. This is compatible
with the convention that $\inf \emptyset = \infty$ which we use
in~\eqref{eq:d-int}.

There is another variant of $\dint$ which measures how far an object
$R$ is from being a retract of another object $X$. More precisely, for
$R, X \in \Ob(\mathscr{C})$ we define
\begin{equation} \label{eq:d-rint} \dret(R,X) = \inf \Bigl\{r \geq 0
  \bigm| \exists \, \varphi: \Sigma^rR \longrightarrow X, \, \exists
  \, \psi: \Sigma^rX
  \longrightarrow R \\
  \;\; \text{such that} \; \psi \circ \Sigma^r \varphi = \eta_{2r}^R
  \Bigr\}.
\end{equation}
In contrast to $\dint$, the measurement $\dret(-,-)$ is in general not
symmetric (such a structure is sometimes called a quasi-pseudometric).

The pseudo-metric $\dint$, admits a shift-invariant version $\sdint$ that will play 
an important role further on:
\begin{equation}\label{eq:sdint}
\sdint (X,Y)=\inf_{r,s} \dint (\Sigma^{r}X,\Sigma^{s}Y)
\end{equation}
It is easily seen that this is also a pseudo-metric. A similar
construction applies also to $\dret$ and produces $\sdret$.


Here is an additional notation to be used later in the paper. Let $Z$
be a set endowed with a (not necessarily symmetric) pseudo-metric $d$
and $X, Y \subset Z$. We denote by

$$d(X,Y) := \sup_{x \in X} \inf_{y \in Y} d(x,y)$$
%$$\vec{d}(X,Y) := \sup_{x \in X} \inf_{y \in Y} d(x,y)$$ 
the longest distance
between the points of $X$ and those of $Y$. 
%(The arrow above the $d$ indicates that
Note that this is not symmetric in $(X,Y)$.
We will sometimes use this measurement for $d=\dint$  and $d=\dret$.
%and denote it $\vec{d}_{\text{int}}$. Similarly
%we can also define also $\vec{d}_{\text{r-int}}$ (although $\dret$ is
%not a symmetric pseudo-metric to start with).


To end this subsection, we list a few basic metric properties of PC
functors in the following lemma which is immediate to establish. In
the statement the same notation $\dint$ is used for the inteleaving
distance in various PC categories.
\begin{lem} \label{l:d-func} Let $\msc', \msc''$ be PC's and
  $\msf: \msc' \longrightarrow \msc''$ a PC-functor.
  \begin{enumerate}
  \item For every $A, B \in \Ob(\msc')$ we have
    $\dint(\msf A, \msf B) \leq \dint(A,B)$.
  \item If there exists a PC functor
    $\msg: \msc'' \longrightarrow \msc'$ such that
    $\dint(\msg \circ \msf, \id_{\msc'}) \leq s$ then for every
    $A, B \in \Ob(\msc')$ we have
    $|\dint(\msf A, \msf B) - \dint (A,B)| \leq 2s$.
  \item If there is a PC-functor $\msh: \msc'' \longrightarrow \msc'$
    such that $\dint(\msf \circ \msh, \id_{\msc''}) \leq r$ then
    $$\dint(\msc'', \textnormal{image\,}\msf) \leq r.$$
  \end{enumerate}
\end{lem}




\subsection{Triangulated persistence categories and TPC
  approximability}\label{subsec:TPC_approx}
We discuss here approximability in the setting of triangulated
persistence categories.

\subsubsection{TPCs} \label{sb:tpc} A {\em Triangulated Persistence
  Category} (TPC in short) $\mathscr{C}$ is a PC with an additional
structure which makes its $0$-level category $\mathscr{C}^0$ a
triangulated category. The precise definition of a TPC involves
several other axioms and we refer the reader to~\cite{BCZ:tpc} for a
detailed exposition of the subject.

Let $\msc$ be a TPC and $S \subset \Ob(\msc)$. We write
$\langle S \rangle^{\Delta}$ for the minimal full sub-TPC of $\msc$
which contains all the objects of $S$ and is closed under
$0$-isomorphisms. (Note that in particular the shift and translation
functors of $\msc$ restrict to respective functors on
$\langle S \rangle^{\Delta}$.) 

%We denote by
%$(\langle S \rangle^{\Delta})^0$ the $0$-level category of the TPC
%$\langle S \rangle^{\Delta}$.

\subsubsection{TPC approximability} \label{subsubsec:approx_TPC} The
key approximability definition for the paper is the following.
\begin{dfn}\label{def:TPC-approx} A pseudo-metric space $(X,d)$ is
  approximable through triangulated persistence categories (in short,
  TPC-approximable) if there exists a constant $A\geq 1$ such that for
  each $\epsilon >0$ the following structure exists. For each
  $0< \eta <\epsilon$ there is a TPC, $\mathscr{Y}_{\epsilon,\eta}$, with two
  properties:
  \begin{itemize}
  \item[i.] There exists an $(A,\eta)$-quasi-isometric embedding (see Remark \ref{rem:def-TPC-approx}):
    $$\Phi_{\epsilon,\eta}:(X,d) \to (\Ob(\mathscr{Y}_{\epsilon,\eta}), \sdint^{\ \epsilon,\eta})$$
  \item[ii.] There exists a {\em finite} family
    $\mathcal{F}_{\epsilon,\eta} \subset \Ob(\mathscr{Y}_{\epsilon,\eta})$ with the
    property:
    $$\sdint^{\ \epsilon,\eta}(\Phi_{\eta}(x), \Ob( \langle\mathcal{F}_{\epsilon,\eta}
    \rangle ^{\Delta}) ) < \epsilon, \ \forall x\in X$$
  \end{itemize} 
  where $\sdint^{\ \epsilon,\eta}$ is the shift invariant interleaving pseudo-metric associated to
  $\mathscr{Y}_{\epsilon,\eta}$ (see \eqref{eq:sdint}).
\end{dfn}

For fixed
$\epsilon$, the data $(\Phi, \F)$ with  $\Phi= \{\Phi_{\epsilon,\eta}\}_{0<\eta <\epsilon}$,
$\F=\{\F_{\epsilon,\eta} \}_{0<\eta<\epsilon}$ is called 
TPC $\epsilon$-{\em approximating data} for $(X,d)$. Thus, $(X,d)$ is
TPC-approximable if $\epsilon$-TPC-approximating data for $(X,d)$
exists for each $\epsilon>0$. It can certainly happen that multiple
choices of such data exist for the same $(X,d)$. When one of the two parameters $\epsilon$ or $\eta$ are clear from the context we omit it from the notation. For instance, for fixed $\epsilon$, we write $\Phi_{\eta}$ for the relevant quasi-isometric embeddings , $\sdint^{\ \eta}$ for the respective pseudo-metrics and so forth. 


\

This definition  is the TPC  version of Definition \ref{def:approx}. 
We will also  use in the paper the notion of  {\em TPC retract approximability}.  TPC retract approximability is defined just as above
without any modifications for point i. above but by using $ \sdret$ at
point ii.  (we emphasize that the {\em only} change is at point ii. in
Definition \ref{def:TPC-approx}). Pairs $(\Phi, \F)$ as above but ensuring retract approximability (in other words, using $\sdret$ instead of $\sdint$ for condition ii in the definition) will be called TPC {\em retract $\epsilon$-approximating data}.

\

Point c. in Remark \ref{rem:relation-TPCapp} below explains more precisely how  TPC  approximability, as just defined, fits with the notion of categorical metric approximability introduced in Definition \ref{def:approx}.


\begin{rem}\label{rem:def-TPC-approx}
  a. Recall that an $(A,B)$-quasi-isometric embedding
  $$(X,d) \to (Y,d')$$ between two metric spaces is a
  map $\Phi : X\to Y$ such that:
  $$\frac{1}{A}\  d(x,y)-B \leq d'(\Phi(x),\Phi(y))\leq A\ d(x,y)+B~.~$$
  Thus, the first point of the Definition \ref{def:TPC-approx} shows that, when $\epsilon$ is
  fixed, the metrics
  $\sdint^{\ \eta}$ defined on the objects of $\mathscr{Y}_{\eta}$, when
  restricted to $X$, get closer and closer to a (pseudo)-metric
  equivalent to $d$, when $\eta\to 0$. More precisely, we can put for
  each $x,y\in X$
  $$\hat{d}(x,y)=\limsup_{\eta\to 0} \sdint^{\ \eta} (x,y)~.~$$
  This gives a pseudo-metric defined on $X$ and point i. of the
  Definition implies that this $\hat{d}$ is equivalent to $d$, the
  original metric on $X$.

  b. A $(A,\eta)$-quasi-isometric embedding is also a
  $(A,\eta')$-quasi-isometric embedding whenever $\eta'\geq \eta$. As
  a result, to verify the property in the definition, one only needs
  to show the existence of $\mathscr{Y}_{\eta}$, $\Phi_{\eta}$, and of
  the families $\F_{\epsilon} =\{\F_{\epsilon,\eta}\}$ whenever $\eta$ is small enough.

  c. This definition will be applied to the case when $(X,d)$ is a
  space of Lagrangian submanifolds endowed with the spectral metric
  $d=d_{\gamma}$ and the categories $\mathscr{Y}_{\eta}$ are TPC
  refinements of derived Fukaya categories (we assume again $\epsilon$ fixed here). 
  The role of $\eta$ is to
  keep track of the size of the perturbations needed to define these
  Fukaya categories. While this parameter is irrelevant in the
  non-filtered case, for purposes of approximations it is crucial to
  keep track of it because any choice of perturbations needed to
  define the Fukaya category determines some $\eta$ which constraints what the approximating accuracy $\epsilon$ can be. Controlling systematically the various TPC refinements of  Fukaya categories  relative to the choices of perturbation data of varying size requires significant elaboration which is contained in \S\ref{s:sys-tpc}. In particular, the metric $\hat{d}$ above is a version  of $\widehat{d}_{\mathrm{int}}$ from Lemma \ref{l:lim-dint}. In this setting of families of filtered $A_{\infty}$-categories it is operationally useful to use a more refined and precise variant of the definition above. This is formulated  in Definition \ref{d:approx-sys} and its retract analogue is  in Definition \ref{d:approx-sys-ret}. See also Remark \ref{rem:different_approx} for a comparison between these notions.
  
  d. The constant $A$
  will be equal to $2$ in the Lagrangian topology application. The  application to Lagrangian topology was determinant for our choice of the shift invariant
 metric $\sdint$ in Definition \ref{def:TPC-approx} (as opposed to
  $\dint$ which could as well have been used, in principle). The
  reason is that the spectral pseudo-metric is shift invariant and
  thus the limit $\hat{d}$ from point a above needs to also be shift
  invariant (see also Remark \ref{rem:relation-TPCapp} ).

  e. In certain contexts, when choices of perturbations are not
  required, and $\epsilon$ is fixed, one presumably can replace the family $\mathscr{Y}_{\eta}$ by a single $PC$ (or TPC) $\mathscr{Y}_{ 0}$ - in other words include
  $\eta=0$ - in the definition above.

  f. In Definition \ref{def:TPC-approx} we have not imposed any
  relation tying the categories $\mathscr{Y}_{\epsilon,\eta}$, when
  $\eta \to 0$, nor did we assume any relation among the respective
  families $\F_{\epsilon,\eta}$.  In practice, our constructions in
  the case of Fukaya categories produce categories that satisfy a
  property more precise than the one in Definition
  \ref{def:TPC-approx}, in the sense that, for fixed $\epsilon >0$, the
  categories $\mathscr{Y}_{\epsilon,\eta}$, with $\eta\to 0$, are related by
  certain comparison functors and the corresponding families
  $\F_{\epsilon,\eta}$ correspond one to the other through these
  functors. This more refined property is more technical to state and
  is formulated in Definition \ref{d:approx-sys}. Further below we
  will also discuss the dependence of our constructions relative to $\epsilon$.
  
\end{rem}

We will also make use of a simplified notion of TPC-approximability
that reformulates point ii of Definition \ref{def:TPC-approx}.


\begin{dfn} \label{def:simple_approx} Fix a triangulated persistence
  category $\msc$ and $\mathcal{X}\subset \Ob(\msc)$.  For $\epsilon >0$, we say that $\mathcal{X}$ is $\epsilon$-{\em approximable in} $\msc$ if there
  exists a {finite} family $\F_{\epsilon}\subset \Ob(\msc)$, called an
  $\epsilon$-{\em approximating family} for $\mathcal{X}$ in $\msc$, such that
  $$\dint\left( x, \Ob (\langle \F_{\epsilon}\rangle^{\Delta})\right)
  < \epsilon, \forall x \in \mathcal{X}.~$$ The set $\mathcal{X}$ is called approximable
  in $\msc$ if it is $\epsilon$-approximable for each $\epsilon >0$.
 We say that $\mathcal{X}$ is  {\em retract $\epsilon$-approximable in}
 $\msc$ if the same property as above holds but with $\dret$ in the place of $\dint$. 
  \end{dfn} 

\begin{rem}\label{rem:relation-TPCapp} a. Recall from \S\ref{sbsb:pc-metrics} that 
both $\dint$ and $\dret$ have
  shift-invariant versions.  Thus, one could use these shift invariant metrics,  $\sdint$ and $\sdret$, instead of $\dint$ and, respectively, $\dret$ in the definition above. However, this has no impact on the notion defined. Indeed, the set $\Ob (\langle \mathcal{S}\rangle ^{\Delta})$ where
  $\langle \mathcal{S} \rangle ^{\Delta}$ is the smallest sub-TPC of
  $\C$ containing $\mathcal{S}\subset \Ob (\C)$ is shift and
  translation invariant. Thus,
  $\dint(x, \Ob (\langle \F_{\epsilon}\rangle^{\Delta})) =\sdint(x,
  \Ob (\langle \mathcal{F}_{\epsilon}\rangle^{\Delta}))$ for all
  $x\in Y$, and similarly for $\dret$.  As a result, the regular and shift
  invariant versions of approximability in the sense of Definition
  \ref{def:simple_approx} coincide and the same is true for retract
  approximability. 
  
  b. Notice that, with the terminology in Definition \ref{def:simple_approx}, point ii in Definition \ref{def:TPC-approx} is equivalent to the fact that $\mathcal{X}=\Phi_{\eta}(X)$ is $\epsilon$-approximable in $\mathscr{Y}_{\eta}$ and similarly for retract approximability ($\epsilon$ is fixed here). 
  
  
  c. The notion of $\epsilon$-approximability
 in Definition \ref{def:simple_approx} obviously implies metric categorical
 $\epsilon$-approximability of $\mathcal{X}$ in $\msc$, as  introduced in Definition \ref{def:approx} (the difference between the two definitions is that in  \ref{def:simple_approx} the metric on the objects of the category $\msc$ is fixed to be the interleaving metric associated with the persistence structure on $\msc$).
 
 d. Assume, as above, that $\msc$ is a TPC.  
  The following statement is an exercise in manipulating exact triangles in the triangulated category $\msc^{0}$:  
  \begin{equation}\label{eq:dint-dret}
  \mathrm{For\ }\  x, y\in \Ob(\msc)\ \mathrm{if}\ \sdret (x,y) <\epsilon \ , \mathrm{then}\ \exists\ k_{x}\in\Ob(\msc),\ \mathrm{such\ that}\  \sdint(x\oplus k_{x},y) < 2\epsilon ~.~
  \end{equation}
It is immediate to see that we also have for all $x,y, k\in \Ob(\msc)$ , $\sdret(x,y)\leq \sdint (x\oplus k,y)$. From (\ref{eq:dint-dret}) we deduce that if $\mathcal{X}$ is retract $\epsilon$-approximable in $\msc$, in the sense of Definition \ref{def:simple_approx}, then it also is retract categorically $\epsilon$-approximable in the sense of \S\ref{subsec:overview}.  In other words, if $\mathcal{X}$ is retract 
$\epsilon$-approximable in $\msc$, then for each $x\in \mathcal{X}$ there exist $k_{x}\in \msc$ and $y\in \Ob (\langle \mathcal{F}_{\epsilon}\rangle^{\Delta}))$ such that $\dint (x\oplus k_{x},y) < 2\epsilon$.

e. Idempotent completion in TPCs is a subtle matter, much more so than its counterpart in triangulated categories. This topic is studied by Miller in \cite{Miller:idempotents} but we will not appeal in this paper to the results there.

 f. In \cite{BCZ:tpc} were introduced a class of so-called fragmentation pseudo-metrics $d^{\F}(-,-)$ on the objects of a TPC that depend on the choice of a family $\F\subset \Ob(\C)$. These pseudo-metrics are closely related to approximability.  Given $\mathscr{C}$ as above, $d^{\F_{\epsilon}}(0,Y)\leq \epsilon/4$ implies that  $\F_{\epsilon}$ is $\epsilon$-approximating. Conversely, if $\F_{\epsilon}$ is $\epsilon$-approximating, then $d^{\F_{\epsilon}}(0,Y)\leq \epsilon$.  Understanding the behaviour of these fragmentation pseudo-metrics when the family $\mathcal{F}$ changes was one of the main motivations leading to the definition of approximability in this paper. \end{rem}
  
 
